\documentclass{beamer}

%% Tema y configuraciones
\usetheme{Boadilla}
\usecolortheme{default}
\setbeamertemplate{navigation symbols}{}
\usepackage{xcolor}
\usepackage{tikz}
\usetikzlibrary{arrows.meta, positioning}
\hypersetup{
    colorlinks=true,
    linkcolor=.,      % color de enlaces internos (índices, refs)
    urlcolor=blue,    % color de URLs
    citecolor=.,      % color de referencias bibliográficas
}

%% Helpers (macros y notación)
\newcommand{\E}{\mathbb{E}}
\newcommand{\Var}{\mathrm{Var}}
\newcommand{\Cov}{\mathrm{Cov}}
\newcommand{\1}{\mathbb{I}}
\newcommand{\MSE}{\mathrm{MSE}}
\DeclareMathOperator*{\argmin}{arg\,min}
\DeclareMathOperator*{\argmax}{arg\,max}

%% Encabezado
\title[BDML]{Big Data y Machine Learning \\ para Economía Aplicada}
\subtitle{Week 07: Aprendizaje No Supervisado y Reducción de Dimensión}
\author{Eduard F. Martinez Gonzalez, Ph.D.}
\institute[]{Departamento de Economía, Universidad Icesi}
\date{\today}

%%=================================================
%% Begin document
\begin{document}

%%=================================================
%% Primer slide
\begin{frame}
\titlepage
\end{frame}

%%=================================================
\begin{frame}{Previously\ldots}
\small
\begin{itemize}\itemsep5pt
  \item Seis sesiones de aprendizaje \textbf{supervisado}: siempre hubo una $y$ que predecir (o un $\alpha$, $\tau(x)$ que estimar) y un error \emph{out-of-sample} contra el cual rendir cuentas.
  \item Con DML y \emph{causal forests} cerramos el bloque predicción--causalidad.
\end{itemize}
\pause

\vspace{0.2cm}
\begin{block}{Hoy: soltamos la $y$}
Solo tenemos los $x$. Dos familias de preguntas:
\begin{enumerate}\itemsep2pt
  \item \textbf{Clustering:} ¿hay \emph{grupos} naturales? (segmentar consumidores, tipologías de municipios)
  \item \textbf{Reducción de dimensión:} ¿hay \emph{pocas dimensiones} que resuman muchas variables? (índices socioeconómicos, factores latentes)
\end{enumerate}
\end{block}
\pause

\vspace{0.1cm}
Y una advertencia desde ya: sin $y$ \textbf{no hay error de prueba} que arbitre. La validación cambia de naturaleza --- y el riesgo de \emph{encontrar} estructura donde no la hay se dispara.
\end{frame}


%%#################################################
%%#################################################
%% PARTE 1: EL PROBLEMA SIN Y
%%#################################################
%%#################################################
\section{El problema no supervisado}
\begin{frame}[noframenumbering]
\tableofcontents[currentsection]
\end{frame}

%%=================================================
\begin{frame}{Qué cambia cuando no hay $y$}
\small
\begin{center}
\footnotesize
\setlength{\tabcolsep}{4pt}
\begin{tabular}{lll}
\hline
 & Supervisado (S1--S6) & No supervisado (hoy) \\
\hline
Objetivo & predecir $y$ / estimar efectos & describir estructura de los $x$ \\
Criterio & riesgo $R(f)$, error de prueba & ¿? (no hay respuesta correcta) \\
Validación & CV, conjunto de prueba & estabilidad, interpretación, uso posterior \\
Peligro & sobreajustar & \textbf{ver patrones donde no los hay} \\
\hline
\end{tabular}
\end{center}
\pause

\vspace{0.25cm}
\textbf{¿Cómo se valida entonces?} Tres sustitutos imperfectos del error de prueba:
\begin{enumerate}\itemsep4pt
  \item \textbf{Estabilidad:} ¿los grupos/componentes sobreviven a remuestrear los datos, cambiar semillas, perturbar variables?
  \item \textbf{Interpretabilidad:} ¿la estructura tiene lectura económica coherente?
  \item \textbf{Uso posterior} (\emph{downstream}): si el índice/cluster luego \emph{predice} algo (pobreza, impago), esa tarea supervisada lo audita.
\end{enumerate}
\pause

\begin{block}{Regla de honestidad}
Los algoritmos de hoy \textbf{siempre} devuelven clusters y componentes --- haya o no estructura real. La carga de la prueba es nuestra.
\end{block}
\end{frame}


%%#################################################
%%#################################################
%% PARTE 2: CLUSTERING
%%#################################################
%%#################################################
\section{Clustering: $k$-means, jerárquico y mezclas}
\begin{frame}[noframenumbering]
\tableofcontents[currentsection]
\end{frame}

%%=================================================
\begin{frame}{$k$-means: el problema de optimización}
\small
\textbf{Objetivo:} partir las $n$ observaciones en $K$ grupos $C_1, \ldots, C_K$ minimizando la dispersión \emph{dentro} de cada grupo:
\[ \min_{C_1, \ldots, C_K} \; \sum_{k=1}^{K} \; \sum_{i \in C_k} \big\| x_i - \mu_k \big\|^2, \qquad \mu_k = \text{centroide (media) del grupo } k \]
\pause

\begin{itemize}\itemsep5pt
  \item Equivalente a minimizar las distancias cuadráticas \emph{intra}-cluster (WSS): grupos compactos y separados.
  \item El problema exacto es \textbf{combinatorio} (NP-duro): hay $\sim K^n$ particiones. Otra vez la salida es una \textbf{heurística} inteligente. \pause
\end{itemize}

\vspace{0.15cm}
\textbf{Algoritmo de Lloyd (1957):}
\begin{enumerate}\itemsep3pt
  \item Inicializar $K$ centroides (aleatorios o \emph{k-means++}).
  \item Repetir hasta converger: (a) \textbf{asignar} cada $i$ a su centroide más cercano; (b) \textbf{recalcular} cada centroide como la media de sus asignados.
\end{enumerate}
\pause

\begin{itemize}\itemsep3pt
  \item Cada paso baja el objetivo $\Rightarrow$ converge\ldots\ a un \textbf{óptimo local}. En la práctica: correr con 20--50 inicios distintos y quedarse con el mejor.
  \item Y como todo lo basado en distancias: \textbf{estandarizar primero} (sesión 2, otra vez).
\end{itemize}
\end{frame}

%%=================================================
\begin{frame}{¿Cuántos grupos? El codo y la silueta}
\small
$K$ no lo dicen los datos solos: el WSS \emph{siempre} baja con $K$ (con $K = n$, cero).

\vspace{0.15cm}
\begin{center}
\begin{tikzpicture}[scale=0.7]
  \draw[->] (0,0) -- (8.6,0) node[below, font=\scriptsize] {$K$};
  \draw[->] (0,0) -- (0,4.6) node[above, font=\scriptsize] {WSS};
  \draw[blue, very thick] plot coordinates {(1,4.2) (2,2.6) (3,1.7) (4,1.45) (5,1.3) (6,1.2) (7,1.12) (8,1.06)};
  \foreach \px/\py in {1/4.2, 2/2.6, 3/1.7, 4/1.45, 5/1.3, 6/1.2, 7/1.12, 8/1.06} \fill[blue] (\px,\py) circle (2.4pt);
  \draw[red, thick] (3,1.7) circle (13pt);
  \node[font=\scriptsize, red] at (3.9,2.5) {el ``codo''};
  \foreach \px in {1,...,8} \node[font=\tiny] at (\px,-0.35) {\px};
\end{tikzpicture}
\end{center}
\pause

\begin{itemize}\itemsep3pt
  \item \textbf{Codo:} elegir el $K$ donde la ganancia marginal se desploma. (¿Les recuerda el \emph{scree plot} que viene en PCA? No es coincidencia.)
  \item \textbf{Silueta} (Rousseeuw, 1987): comparar la distancia media \emph{intra}-cluster de cada $i$ con la del cluster vecino; el promedio $\in [-1,1]$ da una nota por $K$.
  \item Y el criterio que manda: \textbf{el uso} --- si solo se pueden gestionar 4 segmentos, $K = 4$. $K$ es una \emph{decisión}, no un parámetro verdadero.
\end{itemize}
\end{frame}

%%=================================================
\begin{frame}{Clustering jerárquico: el árbol de las fusiones}
\small
\textbf{Aglomerativo:} empezar con $n$ clusters de una observación; en cada paso, \textbf{fusionar los dos más cercanos}; registrar la historia completa.

\vspace{0.15cm}
\begin{center}
\begin{tikzpicture}[scale=0.62]
  % hojas
  \foreach \x/\lbl in {1/A, 2/B, 3.2/C, 4.6/D, 5.6/E} \node[font=\scriptsize] at (\x,-0.4) {\lbl};
  % fusiones: (A,B) h=1; (D,E) h=1.4; (AB,C) h=2.4; (ABC,DE) h=3.4
  \draw[thick] (1,0) -- (1,1) -- (2,1) -- (2,0);
  \draw[thick] (4.6,0) -- (4.6,1.4) -- (5.6,1.4) -- (5.6,0);
  \draw[thick] (1.5,1) -- (1.5,2.4) -- (3.2,2.4) -- (3.2,0);
  \draw[thick] (2.35,2.4) -- (2.35,3.4) -- (5.1,3.4) -- (5.1,1.4);
  % eje de altura
  \draw[->] (7.0,0) -- (7.0,4.0) node[above, font=\scriptsize] {distancia de fusión};
  % corte
  \draw[dashed, red, thick] (0.4,2.9) -- (6.3,2.9);
  \node[font=\scriptsize, red] at (8.0,2.9) {corte: 2 grupos};
\end{tikzpicture}
\end{center}
\pause

\begin{itemize}\itemsep4pt
  \item El \textbf{dendrograma} muestra \emph{todas} las resoluciones a la vez: cortar a una altura $=$ elegir el número de grupos \emph{después} de ver la estructura.
  \item ¿``Cercanos'' entre grupos? El \textbf{linkage}: completo (máx.), promedio, simple (mín.\ --- encadena), Ward (varianza). Completo/promedio/Ward son los sensatos por defecto.
  \item Costo $O(n^2)$ en distancias: para $n$ muy grande, $k$-means; para cientos/miles y estructura anidada (municipios, sectores), jerárquico brilla.
\end{itemize}
\end{frame}

%%=================================================
\begin{frame}{Mezclas de gaussianas y EM, en una diapositiva}
\small
\textbf{La versión probabilística del clustering:} los datos vienen de una \textbf{mezcla} de $K$ normales,
\[ f(x) \;=\; \sum_{k=1}^{K} \pi_k \, \phi\big(x;\, \mu_k, \Sigma_k\big) \]
y pertenecer a un grupo es una \textbf{variable latente}.
\pause

\begin{itemize}\itemsep5pt
  \item \textbf{Asignaciones suaves:} en vez de ``$i$ está en el grupo 2'', $P(\text{grupo } k \mid x_i)$ --- responsabilidades. $k$-means es el caso límite (varianzas iguales $\to 0$: asignación dura). \pause
  \item \textbf{Algoritmo EM} (esbozo): alternar \textbf{E} (calcular responsabilidades dados los parámetros) y \textbf{M} (re-estimar $\pi_k, \mu_k, \Sigma_k$ ponderando por responsabilidades) --- la misma coreografía de Lloyd, en versión probabilística; sube la verosimilitud en cada paso, converge a óptimo local. \pause
  \item \textbf{Qué compra}: clusters elípticos (no solo esféricos), tamaños distintos, incertidumbre de pertenencia --- útil cuando la frontera entre segmentos es gris (¿hogar ``vulnerable'' o ``clase media''?).
\end{itemize}
\pause

\begin{block}{Para profundizar}
La derivación completa del EM: ESL, cap.\ 8.5 --- opcional; el examen no la exige.
\end{block}
\end{frame}

%%=================================================
\begin{frame}{Las advertencias del clustering (léanlas dos veces)}
\small
\begin{enumerate}\itemsep6pt
  \item \textbf{La escala manda.} Sin estandarizar, la variable de mayor varianza \emph{define} los grupos (ingreso en pesos aplasta a años de educación). Y la elección de variables ya es media respuesta: clusters de \emph{consumo} $\neq$ clusters de \emph{demografía}. \pause
  \item \textbf{$k$-means busca esferas.} Grupos alargados, anidados o de densidades muy distintas lo confunden; el jerárquico con Ward o el GMM son más flexibles. \pause
  \item \textbf{Atípicos:} un outlier puede secuestrar un centroide entero. Inspeccionar antes; considerar removerlos o usar medoides. \pause
  \item \textbf{Estabilidad como hábito:} repetir con submuestras del 80--90\%; si los grupos bailan, no eran estructura. \pause
\end{enumerate}

\vspace{0.1cm}
\begin{block}{Y la madre de todas}
El algoritmo \textbf{siempre} entrega $K$ grupos --- también sobre ruido puro. Un cluster no es un hallazgo hasta que sobrevive estabilidad $+$ interpretación $+$ (idealmente) validación \emph{downstream}.
\end{block}
\end{frame}


%%#################################################
%%#################################################
%% PARTE 3: PCA
%%#################################################
%%#################################################
\section{PCA: componentes principales}
\begin{frame}[noframenumbering]
\tableofcontents[currentsection]
\end{frame}

%%=================================================
\begin{frame}{La idea: las direcciones que concentran la varianza}
\small
Con $p$ variables correlacionadas, la información \emph{efectiva} vive en menos dimensiones. PCA busca las \textbf{combinaciones lineales} de máxima varianza:

\vspace{0.15cm}
\begin{center}
\begin{tikzpicture}[scale=0.72]
  \draw[->] (-2.6,0) -- (2.9,0) node[right, font=\scriptsize] {$x_1$};
  \draw[->] (0,-1.7) -- (0,2.2) node[above, font=\scriptsize] {$x_2$};
  \foreach \px/\py in {-1.8/-0.9, -1.4/-0.5, -1.1/-0.8, -0.7/-0.2, -0.4/-0.5, 0/0.1, -0.1/-0.3, 0.4/0.4, 0.7/0.1, 1.0/0.7, 1.3/0.4, 1.7/1.0, 1.9/0.7, 0.2/0.5} \fill[gray] (\px,\py) circle (2pt);
  \draw[-{Stealth}, very thick, blue] (0,0) -- (2.2,1.27);
  \node[font=\scriptsize, blue] at (2.55,1.55) {PC1};
  \draw[-{Stealth}, very thick, red] (0,0) -- (-0.55,0.95);
  \node[font=\scriptsize, red] at (-0.95,1.25) {PC2};
\end{tikzpicture}
\end{center}
\pause

\begin{itemize}\itemsep4pt
  \item \textbf{PC1}: la dirección $w$ que maximiza la varianza de la proyección $w'x$.
  \item \textbf{PC2}: la de máxima varianza \emph{ortogonal} a PC1; y así sucesivamente.
  \item Quedarse con las primeras $M \ll p$: comprimir perdiendo lo mínimo. \pause
  \item Requisitos de aseo: \textbf{centrar} siempre; \textbf{estandarizar} casi siempre (si no, la variable de mayor varianza se roba el PC1 --- misma trampa del clustering).
\end{itemize}
\end{frame}

%%=================================================
\begin{frame}{La derivación: PCA es un problema de autovalores}
\small
Sea $\Sigma$ la matriz de covarianza ($p \times p$) de los datos centrados. Buscamos $w$:
\[ \max_{w} \;\; w' \Sigma\, w \qquad \text{sujeto a} \qquad w'w = 1 \]
(la restricción impide inflar la varianza alargando $w$).
\pause

\textbf{1. Lagrangiano} (¡el de siempre!):
\[ \mathcal{L}(w, \lambda) \;=\; w' \Sigma\, w \;-\; \lambda\,(w'w - 1) \]
\pause
\textbf{2. Condición de primer orden:}
\[ \frac{\partial \mathcal{L}}{\partial w} = 2\,\Sigma w - 2\,\lambda w = 0 \qquad \Longrightarrow \qquad \boxed{\; \Sigma\, w \;=\; \lambda\, w \;} \]
¡$w$ debe ser un \textbf{autovector} de $\Sigma$, con autovalor $\lambda$!
\pause

\textbf{3. ¿Cuál autovector?} Evaluando el objetivo: $w'\Sigma w = \lambda\, w'w = \lambda$.
\begin{itemize}\itemsep2pt
  \item La varianza de la proyección \textbf{es} el autovalor $\Rightarrow$ PC1 $=$ autovector del \textbf{mayor} autovalor $\lambda_1$; PC2, el del segundo (la ortogonalidad sale gratis: $\Sigma$ simétrica); etc.
\end{itemize}
\pause

\begin{block}{Para el Taller 2}
``PCA como problema de autovalores'' es exactamente esta derivación. Sépanla contar de memoria.
\end{block}
\end{frame}

%%=================================================
\begin{frame}{¿Cuántos componentes? Varianza explicada}
\small
Cada componente explica una fracción $\lambda_j / \sum_{k} \lambda_k$ de la varianza total.

\vspace{0.15cm}
\begin{center}
\begin{tikzpicture}[scale=0.68]
  \draw[->] (0,0) -- (7.6,0) node[below, font=\scriptsize] {componente $j$};
  \draw[->] (0,0) -- (0,4.2) node[above, font=\scriptsize] {$\lambda_j / \sum_k \lambda_k$};
  \draw[blue, very thick] plot coordinates {(1,3.4) (2,1.6) (3,0.8) (4,0.45) (5,0.3) (6,0.2)};
  \foreach \px/\py in {1/3.4, 2/1.6, 3/0.8, 4/0.45, 5/0.3, 6/0.2} \fill[blue] (\px,\py) circle (2.4pt);
  \draw[red, thick] (2,1.6) circle (12pt);
  \node[font=\scriptsize, red] at (3.3,2.3) {codo del \emph{scree plot}};
  \foreach \px in {1,...,6} \node[font=\tiny] at (\px,-0.35) {\px};
\end{tikzpicture}
\end{center}
\pause

\textbf{Los tres criterios de siempre (ninguno automático):}
\begin{itemize}\itemsep3pt
  \item \textbf{Umbral acumulado:} conservar componentes hasta explicar 80--90\%.
  \item \textbf{El codo} del \emph{scree plot} (el mismo razonamiento del WSS).
  \item \textbf{Interpretabilidad:} un componente que no se puede nombrar rara vez sirve para comunicar.
\end{itemize}
\pause
Y si el destino es un modelo supervisado (regresión sobre componentes): elegir $M$ \textbf{por CV} --- ahí sí hay árbitro.
\end{frame}

%%=================================================
\begin{frame}{El uso económico estrella: índices socioeconómicos}
\small
\textbf{Filmer \& Pritchett (2001):} sin datos de ingreso confiables, construir un índice de riqueza como el \textbf{PC1 de los activos del hogar} (materiales de la vivienda, electrodomésticos, acceso a servicios).
\pause

\begin{itemize}\itemsep5pt
  \item \textbf{Las cargas} (\emph{loadings}) del PC1 son los \emph{pesos} del índice: qué activos ``pesan'' más en la riqueza latente. Leerlas \textbf{es} interpretar el componente.
  \item La lógica: la dimensión de máxima covariación entre activos \emph{es} el nivel socioeconómico --- el factor latente que los mueve a todos. \pause
  \item Es el espíritu de los \emph{proxy means tests} (Sisbén incluido): resumir muchas señales observables en un índice focalizable.
\end{itemize}
\pause

\vspace{0.1cm}
\textbf{Letras menudas del oficio:}
\begin{itemize}\itemsep3pt
  \item El \textbf{signo} de cada componente es arbitrario (autovector por $-1$ sigue siendo autovector): orientarlo para que ``más $=$ más riqueza''.
  \item Con variables en escalas distintas: PCA sobre la matriz de \textbf{correlación} (estandarizar).
  \item El índice se \emph{audita} con validación \emph{downstream}: ¿predice consumo, salud, educación? (la sesión 1 saluda).
\end{itemize}
\end{frame}

%%=================================================
\begin{frame}{PCA como compresión --- y el puente a lo que viene}
\small
\textbf{Tres usos del mismo objeto:}
\begin{enumerate}\itemsep4pt
  \item \textbf{Describir}: 2--3 componentes con nombre (riqueza, urbanidad, formalidad).
  \item \textbf{Visualizar}: proyectar en (PC1, PC2) --- el mapa lineal de los datos.
  \item \textbf{Comprimir para predecir}: regresión sobre componentes (PCR) --- reducir $p$ antes del modelo supervisado; $M$ por CV. Prima hermana de Ridge (ambas encogen en las direcciones de poca varianza).
\end{enumerate}
\pause

\vspace{0.2cm}
\begin{block}{El puente conceptual hacia las sesiones 8--9}
PCA aprende una \textbf{representación de baja dimensión} de los datos --- \emph{lineal}. Las redes neuronales aprenderán representaciones \textbf{no lineales} (un \emph{autoencoder} es ``PCA no lineal''), y los \emph{embeddings} de palabras de la sesión 9 son exactamente eso: coordenadas aprendidas en un espacio de baja dimensión donde la geometría significa algo. Guarden esta frase: \textbf{representar es comprimir con sentido}.
\end{block}
\end{frame}

%%=================================================
\begin{frame}{t-SNE y UMAP: mapas no lineales (solo para los ojos)}
\small
Cuando la estructura no es lineal (clusters curvos, variedades), PCA la aplana mal. \textbf{t-SNE} (van der Maaten \& Hinton, 2008) y \textbf{UMAP} (McInnes et al., 2018) construyen un mapa 2D que preserva \textbf{vecindades locales}: quién está cerca de quién.
\pause

\begin{itemize}\itemsep5pt
  \item \textbf{Para qué sirven:} exploración visual --- ver de un vistazo si hay grupos, subgrupos, transiciones. Espectaculares con datos de alta dimensión (imágenes, texto, genómica). \pause
  \item \textbf{Para qué NO sirven} (y aquí se cometen papers enteros de errores):
  \begin{itemize}\itemsep2pt
    \item Las \textbf{distancias globales} en el mapa no significan nada (dos islas lejanas no son ``muy distintas'' por estar lejos).
    \item El \textbf{tamaño y densidad} de los grupos en el mapa es artefacto de los hiperparámetros (\emph{perplexity} / vecinos).
    \item Corridas distintas dan mapas distintos: no es un espacio con coordenadas interpretables.
  \end{itemize} \pause
\end{itemize}

\begin{block}{Regla práctica}
UMAP/t-SNE para \textbf{mirar}; $k$-means/jerárquico sobre los datos (o sobre componentes de PCA) para \textbf{agrupar}; PCA para \textbf{construir índices}. Cada herramienta en su puesto.
\end{block}
\end{frame}


%%#################################################
%%#################################################
%% PARTE 4: SÍNTESIS
%%#################################################
%%#################################################
\section{Síntesis}
\begin{frame}[noframenumbering]
\tableofcontents[currentsection]
\end{frame}

%%=================================================
\begin{frame}{Recap}
\small
\begin{enumerate}\itemsep6pt
  \item Sin $y$ no hay árbitro \emph{out-of-sample}: la validación es \textbf{estabilidad $+$ interpretación $+$ uso posterior} --- y el algoritmo siempre ``encuentra'' algo.
  \item \textbf{$k$-means}: minimizar WSS con Lloyd (óptimos locales $\Rightarrow$ muchos inicios); $K$ por codo/silueta \emph{y} por uso; estandarizar siempre.
  \item \textbf{Jerárquico}: el dendrograma da todas las resoluciones; el linkage importa; cortar es decidir.
  \item \textbf{GMM/EM}: clustering probabilístico con pertenencias suaves; $k$-means es su caso límite.
  \item \textbf{PCA} $=$ autovalores de $\Sigma$ (sépanlo derivar): la varianza de cada componente es su $\lambda$; componentes por codo/umbral/CV; el PC1 de activos es el índice socioeconómico clásico (Filmer--Pritchett).
  \item \textbf{t-SNE/UMAP}: mapas locales para explorar --- jamás para medir distancias ni tamaños.
  \item El hilo a futuro: \textbf{representar es comprimir con sentido} --- PCA lineal hoy; redes y \emph{embeddings} no lineales en las sesiones 8--9.
\end{enumerate}
\end{frame}

%%=================================================
\begin{frame}
\vspace{2.0cm}
\Large{Preparación para la próxima clase\ldots}
\end{frame}

%%#################################################
%%#################################################
%% PARTE 5: PRÓXIMA CLASE
%%#################################################
%%#################################################

%%=================================================
\begin{frame}{Para aprovechar al máximo la próxima sesión}
\small
\textbf{Lecturas obligatorias de esta semana:}
\begin{itemize}\itemsep4pt
  \item ISLR, cap.\ 12 (secs.\ 12.1, 12.2 y 12.4). \\ \textcolor{gray}{Guía: PCA y clustering con las figuras que vimos hoy.}
  \item Filmer \& Pritchett (2001). \\ \textcolor{gray}{Guía: secciones 1--3; fíjense cómo \emph{validan} el índice sin observar ingreso.}
\end{itemize}

\textbf{Recomendadas:} Jolliffe (2002), caps.\ 1--2; McInnes et al.\ (2018, UMAP) para curiosos.
\pause

\vspace{0.2cm}
\textbf{Próxima sesión --- empieza el bloque final:} \textbf{redes neuronales}. Del perceptrón a la retropropagación (que derivaremos: va al Taller 2), y de ahí a las redes convolucionales que leen \textbf{imágenes satelitales}. Lectura previa: Goodfellow, Bengio \& Courville, cap.\ 6 (secciones 6.1--6.3, sin pánico).
\pause

\vspace{0.2cm}
\textbf{Aplicación de esta semana en R:} segmentar hogares por patrones de consumo ($k$-means $+$ jerárquico) y construir un índice socioeconómico con PCA, comparándolo con la clasificación oficial.
\end{frame}

%%=================================================
%% Referencias
\begin{frame}{Referencias esenciales}
\footnotesize
\begin{itemize}\itemsep2pt
  \item James, G., Witten, D., Hastie, T., \& Tibshirani, R. (2021). \emph{An Introduction to Statistical Learning} (2.\ ed.). Springer. Cap.\ 12. [ISLR]
  \item Hastie, T., Tibshirani, R., \& Friedman, J. (2009). \emph{The Elements of Statistical Learning} (2.\ ed.). Springer. Cap.\ 14. [ESL]
  \item Jolliffe, I. (2002). \emph{Principal Component Analysis} (2.\ ed.). Springer. Caps.\ 1--2.
  \item Rousseeuw, P. (1987). Silhouettes: A Graphical Aid to the Interpretation and Validation of Cluster Analysis. \emph{J.\ of Computational and Applied Mathematics}, 20, 53--65.
  \item Filmer, D., \& Pritchett, L. (2001). Estimating Wealth Effects Without Expenditure Data --- or Tears. \emph{Demography}, 38(1), 115--132.
  \item van der Maaten, L., \& Hinton, G. (2008). Visualizing Data using t-SNE. \emph{JMLR}, 9, 2579--2605.
  \item McInnes, L., Healy, J., \& Melville, J. (2018). UMAP: Uniform Manifold Approximation and Projection. \emph{arXiv:1802.03426}.
\end{itemize}
\normalsize
\end{frame}

%%=================================================
%% End document
\end{document}
